\section{Related Work}
\textbf{Streaming encoders.} Chunk-based Conformer and Emformer architectures process audio in fixed
windows to bound latency \citep{okonkwo2024conformer}. Memory-augmented variants carry a summary vector
across chunks to approximate long-range context without violating causality
\citep{bergman2018lstm,kowalski2022diarization}.

\textbf{Look-ahead and dual-mode models.} Dual-mode training jointly optimizes a single encoder for both
streaming and offline decoding, sharing parameters across latency regimes \citep{martinez2022rnnt}.
Fixed uniform look-ahead has been explored as a simple latency-accuracy knob
\citep{fischer2020ctc,silva2019vad}, but, to our knowledge, no prior work varies the look-ahead
allocation by depth within the encoder stack, which is the central mechanism of CBCI.

\textbf{Self-supervised and adaptive representations.} Self-supervised pretraining
\citep{tanaka2021selfsup,schmidt2023wav2vec} and few-shot accent adaptation \citep{nakamura2023fewshot}
improve robustness of downstream ASR; these techniques are complementary to CBCI and are combined with
it in our experiments via a pretrained encoder initialization.

\textbf{Prosody and noise robustness.} Prosody modeling for synthesis \citep{gomez2020prosody} and
reverberant augmentation for recognition \citep{andersson2021noise} motivate our choice of evaluation
conditions, described in Section~\ref{sec:experiments}.
